\documentclass[11pt]{article}
\usepackage[utf8]{inputenc}
\usepackage{amsmath,amssymb}
\usepackage{enumitem}
\usepackage{geometry}
\geometry{margin=1in}

% Custom macros for questions
\newcommand{\QBlock}[3]{%
  \subsection*{Question #1}
  #2
  \begin{enumerate}[label=(\Alph*)]
  #3
  \end{enumerate}
  \vspace{1em}
}

\newcommand{\Correct}[2]{%
  \item \textbf{#1} \hfill $\checkmark$ \\
  \textit{Explanation: #2}
}

\newcommand{\Wrong}[2]{%
  \item #1 \\
  \textit{Explanation: #2}
}

\begin{document}

\title{Isomorphic Physics Questions}
\author{Generated by AI}
\date{\today}
\maketitle

% Questions go here using \QBlock, \Correct, \Wrong macros

\QBlock{1}{A dolphin leaps out of the water with an initial speed $u_0$ at an angle $\alpha$ to the horizontal. What is the dolphin's speed at the highest point of its trajectory?}{
  \Wrong{$u_0$}{At maximum height, only the horizontal component of velocity remains; $u_0$ is incorrect as it is the total initial speed.}
  \Correct{$u_0 \cos \alpha$}{Correct because at the highest point, the vertical component of speed is zero, and only the horizontal component ($u_0 \cos \alpha$) remains.}
  \Wrong{$u_0 \sin \alpha$}{Incorrect because $u_0 \sin \alpha$ is the initial vertical component, which becomes zero at maximum height.}
  \Wrong{0}{Incorrect because the horizontal speed is non-zero at the apex.}
}

\QBlock{2}{A space probe travels with an initial velocity of $\vec{u}_0 = (30.5 \text{ m/s})\hat{i}$ when thrusters provide an acceleration of $\vec{b} = (1.5 \text{ m/s}^2)\hat{i}+(5.0 \text{ m/s}^2)\hat{j}$. Which statement about the probe's velocity is true?}{
  \Wrong{The x-component of the probe’s velocity remains constant.}{Incorrect because there is acceleration in the x-direction, so velocity changes.}
  \Wrong{The y-component of the probe’s velocity remains unchanged.}{Incorrect since there is acceleration in the y-direction, altering the velocity.}
  \Correct{Both the speed and the motion direction change.}{Correct because accelerations in both directions affect both speed and direction.}
  \Wrong{None of these statements is correct.}{Incorrect because option C accurately describes the situation.}
}

\end{document}